\documentclass[10pt]{article}
\providecommand{\FIGDIR}{fig}
\newcommand{\DIGESTDATE}{9 August 2026}
% ===================== PM Research Watch - shared preamble =====================
\usepackage[T1]{fontenc}
\usepackage[utf8]{inputenc}
\usepackage{amsmath,amssymb}
\usepackage{tgpagella}
\usepackage{tgheros}
\renewcommand{\sfdefault}{qhv}
\renewcommand{\ttdefault}{lmtt}
\usepackage{microtype}
\usepackage[english]{babel}

\usepackage[a4paper,top=17mm,bottom=17mm,left=17mm,right=17mm,
            headsep=6mm,footskip=10mm]{geometry}
\usepackage{graphicx}
\usepackage{xcolor}
\usepackage{colortbl}
\usepackage{tikz}
\usetikzlibrary{calc}
\usepackage[most]{tcolorbox}
\usepackage{titlesec}
\usepackage{fancyhdr}
\usepackage{booktabs}
\usepackage{array}
\usepackage{longtable}
\usepackage{enumitem}
\usepackage{multicol}
\usepackage{ragged2e}
\usepackage[hidelinks]{hyperref}

% ------------------------------------------------------------------ palette
\definecolor{Ink}    {HTML}{14262E}
\definecolor{Deep}   {HTML}{0F4C5C}
\definecolor{Teal}   {HTML}{1F7A8C}
\definecolor{Sky}    {HTML}{4FA3C4}
\definecolor{Sage}   {HTML}{6FA287}
\definecolor{Amber}  {HTML}{E3A008}
\definecolor{Clay}   {HTML}{B5643C}
\definecolor{Coral}  {HTML}{D1495B}
\definecolor{Violet} {HTML}{6A5B9E}
\definecolor{Slate}  {HTML}{4A6572}
\definecolor{Paper}  {HTML}{FBF9F5}
\definecolor{Mist}   {HTML}{EEEAE2}
\definecolor{Rule}   {HTML}{D9D3C9}

\pagecolor{Paper}
\color{Ink}

\hypersetup{colorlinks=true, linkcolor=Deep, urlcolor=Teal, citecolor=Teal,
            pdftitle={PM Research Watch}, pdfauthor={Automated daily digest}}

% ------------------------------------------------------------------ headings
\titleformat{\section}
  {\sffamily\bfseries\large\color{Deep}}{}{0pt}
  {}[{\vspace{1.2mm}\color{Amber}\titlerule[1.1pt]}]
\titlespacing*{\section}{0pt}{5.5mm}{2.8mm}

\titleformat{\subsection}
  {\sffamily\bfseries\normalsize\color{Teal}}{}{0pt}{}
\titlespacing*{\subsection}{0pt}{3.6mm}{1.4mm}

\setlist[itemize]{leftmargin=4.2mm, itemsep=1.1mm, topsep=1.3mm, parsep=0pt,
                  label=\textcolor{Teal}{\footnotesize$\blacktriangleright$}}

% ------------------------------------------------------------------ header/footer
\pagestyle{fancy}
\fancyhf{}
\renewcommand{\headrulewidth}{0pt}
\renewcommand{\footrulewidth}{0pt}
\fancyhead[L]{\sffamily\scriptsize\color{Slate}PM RESEARCH WATCH\;\textcolor{Rule}{\textbar}\;\DIGESTDATE}
\fancyhead[R]{\sffamily\scriptsize\color{Slate}Particulate matter monitoring \& health effects}
\fancyfoot[C]{\sffamily\scriptsize\color{Slate}\thepage}

% ------------------------------------------------------------------ components
% metric tile
\newcommand{\metric}[3]{%
  \begin{tcolorbox}[enhanced, colback=white, colframe=Rule, boxrule=0.5pt,
    arc=1.6mm, left=2.2mm, right=2.2mm, top=1.8mm, bottom=1.8mm,
    borderline west={1.7mm}{0pt}{#3}, nobeforeafter, valign=center,
    height=17mm, width=\linewidth]
    \RaggedRight
    {\sffamily\bfseries\fontsize{17}{17}\selectfont\color{#3}#1}\\[0.5mm]
    {\sffamily\fontsize{6.5}{7.4}\selectfont\color{Slate}#2}
  \end{tcolorbox}}

% pull-quote / key insight box
\newtcolorbox{keybox}[1][]{enhanced, colback=Mist, colframe=Mist,
  arc=1.6mm, left=3.2mm, right=3.2mm, top=2.4mm, bottom=2.4mm,
  borderline west={1.7mm}{0pt}{Amber}, #1}

% subtopic band
\newcommand{\band}[2]{%
  \vspace{3.2mm}\noindent
  \begin{tcolorbox}[enhanced, colback=#2, colframe=#2, arc=1.2mm,
    left=2.6mm, right=2.6mm, top=1.1mm, bottom=1.1mm, boxrule=0pt,
    width=\linewidth, nobeforeafter]
    {\sffamily\bfseries\small\color{white}#1}
  \end{tcolorbox}\vspace{1.6mm}}

% one paper entry: {short cite}{title}{journal}{doi}{body}
\newcommand{\paperentry}[5]{%
  \par\vspace{1.9mm}\noindent
  {\sffamily\bfseries\footnotesize\color{Ink}#1}\;%
  {\sffamily\fontsize{7}{8}\selectfont\color{Slate}#3}\\[0.4mm]
  {\footnotesize\itshape\color{Deep}#2}\\[0.6mm]
  {\footnotesize #4}\\[0.3mm]
  {\fontsize{6.6}{7.6}\selectfont\color{Slate}\texttt{doi:}\,\href{https://doi.org/#5}{\texttt{\detokenize{#5}}}}
  \par}

% caption style
\newcommand{\figcap}[1]{\par\vspace{1.0mm}
  {\fontsize{7.2}{8.6}\selectfont\color{Slate}\textsf{\textbf{Fig.}}\; #1\par}\vspace{1.5mm}}

\providecommand{\permil}{\textperthousand}

\setlength{\parindent}{0pt}
\setlength{\parskip}{0pt}
\renewcommand{\arraystretch}{1.18}


\begin{document}

% ==================================================================== MASTHEAD
\thispagestyle{empty}
\begin{tikzpicture}[remember picture, overlay]
  \fill[Deep] (current page.north west) rectangle
      ($(current page.north east)+(0,-67mm)$);
  \fill[Amber] ($(current page.north west)+(0,-67mm)$) rectangle
      ($(current page.north east)+(0,-68.4mm)$);
  \foreach \x/\y/\r in {22/12/0.9, 41/26/1.5, 63/9/0.7, 88/31/1.1, 112/17/1.9,
                        138/29/0.8, 159/13/1.3, 178/54/1.0, 196/21/1.6,
                        30/44/0.6, 74/57/1.2, 125/48/0.9, 168/6/0.8, 191/59/0.7,
                        52/18/1.0, 100/60/0.6, 148/20/1.1, 66/49/1.4, 205/45/0.8}
    \fill[Sky, opacity=0.30] ($(current page.north west)+(\x mm,-\y mm)$) circle (\r mm);
\end{tikzpicture}

\vspace*{4mm}
{\sffamily\fontsize{7.6}{9}\selectfont\color{Sky}%
 AUTOMATED DAILY LITERATURE DIGEST\;\textbullet\;PhD RESEARCH WATCH\par}
\vspace{2.2mm}
{\sffamily\bfseries\fontsize{27}{29}\selectfont\color{white}%
 Particulate Matter\par}
{\sffamily\fontsize{27}{29}\selectfont\color{Sky}%
 Monitoring \& Health Effects\par}
\vspace{5.0mm}
{\sffamily\fontsize{8.4}{10}\selectfont\color{white}%
 Issue 2026-08-09\;\;\textcolor{Amber}{\textbar}\;\;Entry window 9 August 2026\;\;%
 \textcolor{Amber}{\textbar}\;\;12 records in scope\par}
\vspace{18mm}

% -------------------------------------------------------------- metric strip
\noindent
\begin{minipage}[t]{0.190\linewidth}\metric{12}{RECORDS IN SCOPE\\(16 REJECTED)}{Deep}\end{minipage}\hfill
\begin{minipage}[t]{0.190\linewidth}\metric{8}{SENSING / EXPOSURE\\RECORDS (67\%)}{Teal}\end{minipage}\hfill
\begin{minipage}[t]{0.190\linewidth}\metric{6}{SUBTOPIC CLUSTERS\\REPRESENTED}{Sky}\end{minipage}\hfill
\begin{minipage}[t]{0.190\linewidth}\metric{3}{TIER-A\\RECORDS}{Amber}\end{minipage}\hfill
\begin{minipage}[t]{0.190\linewidth}\metric{0}{EFFECT ESTIMATES\\EXTRACTED}{Clay}\end{minipage}

\vspace{6mm}

% -------------------------------------------------------------- signal of day
\section*{\sffamily\bfseries\color{Deep}Signal of the day}
\vspace{-1.5mm}{\color{Amber}\rule{\linewidth}{1.1pt}}\vspace{2.5mm}

\begin{keybox}
{\sffamily\bfseries\small\color{Deep}A hardware day. Three of twelve records are new
instruments, and one of them is the first paper in 158 records to build a device whose
stated purpose is to put an aerosol and a biological endpoint in the same
apparatus.}\\[1.4mm]
\footnotesize
The standing gap this watch has logged since 26 July --- restated in the 2--8 August
weekly over 146 records --- is that \textbf{not one paper coupled an instrument to a
health endpoint}. Shirkhani et al.\ (\textit{J Aerosol Sci}) describe MARTHA, a
high-throughput multipattern aerosol deposition device built explicitly ``for respiratory
toxicity and health-risk assessment''; the author list (Amanatidis, Eiguren-Fernandez,
Spielman) is an aerosol-instrumentation group, not a toxicology group. The record carries
\textbf{no abstract in Crossref, Europe PMC or Semantic Scholar}, so it is listed and not
summarised, and the claim above rests on a title and an author list. That is exactly the
right amount of weight to put on it --- but it is the first time the empty cell in the
f4 cross-tabulation has had a candidate at all.
Alongside it, Li et al.\ (\textit{Photoacoustics}) close the split-instrument problem the
other way: two integrating spheres serving simultaneously as scattering cavities and as a
low-frequency differential Helmholtz resonator, with the connecting tube carrying a
high-frequency longitudinal mode, so absorption and scattering at \textbf{450 and
532\,nm} are measured \emph{in the same volume on one sampling path}
(LOD \textbf{1.2} and \textbf{1.32\,Mm$^{-1}$} at 1\,s, 3$\sigma$). Laffolley et al.\
(\textit{ACS Omega}) go the opposite direction on cost: a stereolithography-printed
two-stage virtual impactor, disposable by incineration, CFD-optimised to
\textbf{9.0} and \textbf{1.3\,$\mu$m} cutoffs. Three different answers to the same
question --- what does a size- or wavelength-resolved measurement have to cost?
\end{keybox}

% Proof 1 left p.2 90 % empty: the six-bullet brief cannot fit under the keybox on p.1
% and spilled its last two bullets alone onto p.2. Per the 06/07/08 Aug log, prose
% trimming does not converge on a spill this size; moving the whole block across the
% break does. The brief now owns p.2.
\clearpage

% -------------------------------------------------------------- exec brief
\section*{\sffamily\bfseries\color{Deep}What changed today}
\vspace{-1.5mm}{\color{Amber}\rule{\linewidth}{1.1pt}}\vspace{2.5mm}

\footnotesize
\begin{itemize}
\item \textbf{Zero effect estimates. The forest plot is absent from this issue for the
first time since the watch began.} This is a reporting-format fact, not an absence of
quantitative content. Three records carry intervals --- Urso et al.'s
\textbf{41\,\% (24.4--57.4)} of EU-27 dementia DALYs attributable to modifiable risk
factors, and their \textbf{7.76\,M (6.69--8.72)} and \textbf{12.28\,M (10.44--13.90)}
prevalence counts --- but a population-attributable fraction and a case count are not
ratios and have no defensible position on a log axis whose null is 1.0. The same rule
excluded mmHg mean differences on 8 August and an OR spanning a factor of fifteen on
7 August. Ni et al.'s $\rho=0.28$ and AUC 0.607 are likewise not ratio-scale. Rather than
render one figure dishonestly, \texttt{f5\_forest} is not generated and the numbers are
carried in the digest text and the register.

\item \textbf{Sunday, and the thinnest window on record after yesterday's: 12 records
against 14 on Saturday and 26 on Friday.} The PubMed connector returned \textbf{6 on the
health axis and 4 on the instrumentation axis}, ending three consecutive empty sensing
queries from that leg; a third broad MeSH sweep returned 15 for the whole day, of which
4 were already shipped in the 4 August issue and re-surfaced on re-indexing, and 4 were
off-topic. The local harvester returned 5 and 1. \textbf{Crossref by ISSN across
17 journals returned 9, of which 3 carried} --- a materially lower yield than the
6-of-7 it delivered on Saturday.

\item \textbf{The instrumentation axis carried the issue, and two thirds of it came from
outside PubMed's biomedical index.} Eight of twelve records are measurement-side
(3 sensing, 3 exposure assessment, 2 occupational). The two strongest instrument papers
reached the watch by different routes --- Li et al.\ through PubMed only because
\textit{Photoacoustics} is biomedically indexed, Laffolley et al.\ only through Europe
PMC's \texttt{CREATION\_DATE} leg. Neither would have been found by the health query.

\item \textbf{Three of twelve records still carry no abstract anywhere, and all three are
on the measurement side.} \textit{Atmospheric Environment} (Ray, size-resolved
hygroscopicity and CCN activity), \textit{Journal of Aerosol Science} (Shirkhani, MARTHA)
and \textit{Journal of Hazardous Materials} (Wei, size-resolved oxidative potential) ---
all same-day Elsevier deposits, all returning HTTP 404 from Semantic Scholar. Separately,
\textbf{the six Elsevier DOIs flagged for re-check on 8 August still have no abstract in
Crossref or Europe PMC 24 hours later}, so the deposit lag is longer than a day and the
re-check should move to a weekly sweep rather than a daily one. Three records held from
7--8 August \emph{did} resolve overnight (\texttt{reveh-2026-0032},
\texttt{envpol.2026.128933}, \texttt{acsomega.6c00972}) --- but all three had already
shipped, so the abstracts arrived too late to be used.

\item \textbf{Nothing new on calibration or drift for a second consecutive day, and the
reason is now measured rather than suspected.} Three Consensus sweeps (calibration and
drift, satellite--surface fusion, humidity and hygroscopic-growth correction) returned
20 hits each, capped at 3. All nine returned DOIs resolved to Crossref creation dates
between \textbf{12 January and 15 July 2026} --- none within four weeks of the window,
let alone inside it. Consensus is a relevance index, not a recency index; it cannot see
a single day and should stop being asked to. \textbf{OpenAlex returned HTTP 429 for the
seventh consecutive run.} arXiv returned 50 records, none published on or after 1 August.

\item \textbf{Trial watch: zero registrations or updates.} Both ClinicalTrials.gov
queries (air-pollution/PM conditions; purifier, filtration, HEPA, cookstove and mask
interventions) returned 0 for 9 August. \texttt{state/trials.json} is unchanged at the
21 registrations reported in the 2--8 August weekly.
\end{itemize}

\clearpage

% ==================================================================== ANALYTICS
\section*{\sffamily\bfseries\color{Deep}Corpus analytics}
\vspace{-1.5mm}{\color{Amber}\rule{\linewidth}{1.1pt}}\vspace{3mm}

\noindent\includegraphics[width=\linewidth]{\FIGDIR/f1_subtopics.png}
\figcap{Subtopic assignment of the 12 in-scope records, single-label by dominant
contribution. \textit{Sensing, forecasting \& instrumentation} and \textit{Exposure
assessment \& modelling} tie at 3 and \textit{Occupational \& indoor} and \textit{Burden,
policy \& mitigation} at 2; \textbf{8 of 12 records are on the measurement side, the
highest share this watch has recorded}. Four of the ten canonical clusters are empty
--- \textit{Neuro / mental health}, \textit{Cardiovascular \& metabolic},
\textit{Reproductive \& developmental} and \textit{Other clinical endpoints}. The
dementia-burden record is deliberately filed under \textit{Burden} rather than
\textit{Neuro}, because its contribution is an attributable-fraction estimate, not a
neurological finding.}

\vspace{2mm}
\noindent\includegraphics[width=\linewidth]{\FIGDIR/f2_design_endpoint.png}
\figcap{Left --- study architecture. \textbf{\textit{Metadata only} and \textit{Measurement
campaign} tie at 3 of 12}, with \textit{Chamber / laboratory} at 2 and
\textit{Observational --- cross-sectional} at 2; \textit{Modelling / inventory} and
\textit{Experimental / toxicology} carry 1 each. Six architectures across twelve records
is the widest spread the watch has rendered at this $n$. \textbf{No observational cohort
and no controlled trial entered today}, for the second consecutive issue. Right --- health
endpoint. \textbf{Eight of twelve records carry \emph{no} human health endpoint}: two
instrument-development papers, four pure exposure-monitoring records, one \emph{in vitro}
oxidative-potential assay and one emissions bench study. The four that do are the
respiratory \emph{in vitro} deposition device, the multi-system mining survey, the
neuro/cognitive burden estimate and the respiratory clinical series --- so the endpoint
panel describes four records and is not a field-level distribution.}

\vspace{2mm}
\noindent\includegraphics[width=\linewidth]{\FIGDIR/f3_metric_geography.png}
\figcap{Left --- particle metric, and the panel is \textbf{measuring this watch's taxonomy
rather than the literature}. Ten of twelve records land in \textit{PM (unspeciated) /
emissions} against one \textit{PM$_{2.5}$ only} (Ni et al.) and one \textit{PM$_{2.5}$ +
PM$_{10}$ jointly} (Adeniji et al.). The classifier has no bin for the quantities the
day's instruments actually report --- absorption and scattering coefficient,
aerodynamic-diameter class, size-resolved hygroscopicity ($\kappa$), size-resolved
oxidative potential, gravimetric mass with SEM morphology --- so six physically distinct
metrics collapse into one bar. \textbf{The figure is correct and the vocabulary is not};
this is logged as a defect against \texttt{plots.py}, not read as a finding. Right ---
geographic provenance. Europe and China tie at 3, South Asia 2, and \textbf{Sub-Saharan
Africa appears for the first time in six issues} (Abeokuta, Nigeria; zero in each issue
from 4 to 8 August). North America contributes 1, its lowest share of the watch.}

\vspace{2mm}
\noindent\includegraphics[width=\linewidth]{\FIGDIR/f6_lifecourse.png}
\figcap{Life-course windows addressed. Counts are non-exclusive and total 4, not 12,
because eight records address no human life stage at all. \textbf{The three perinatal and
paediatric stages are empty for the first time in this watch}; at $n=12$ on a Sunday this
is sampling, not a field-level absence, and the reproductive/developmental literature was
present as recently as 8 August.}

\vspace{2mm}
\noindent\includegraphics[width=\linewidth]{\FIGDIR/f4_heatmap.png}
\figcap{Subtopic $\times$ study-architecture cross-tabulation. Two features are worth
naming. First, the \textit{Metadata only} column is again distributed across
\textit{Sensing}, \textit{Exposure assessment} and \textit{Mechanistic toxicology} rather
than across the health rows --- the fifth consecutive issue in which the least
characterisable records are the most methodologically relevant ones. Second, the
\textit{Sensing} row is populated only in \textit{Metadata only} (1) and \textit{Chamber /
laboratory} (2) --- every sensing record today is either an uncharacterisable deposit or a
bench instrument, and \textbf{not one is a field co-location or a deployment}, which is
the architecture a distributed network is actually built from. The instrument
$\times$ health-endpoint coupling the 2--8 August weekly identified as structurally empty
is still not visible here, because MARTHA is filed under \textit{Metadata only} and cannot
be scored on architecture.}

\clearpage

% ==================================================================== DIGEST
\section*{\sffamily\bfseries\color{Deep}Paper-level digest}
\vspace{-1.5mm}{\color{Amber}\rule{\linewidth}{1.1pt}}\vspace{1mm}

{\fontsize{7.4}{9}\selectfont\color{Slate}Ordered by cluster. Each entry gives the
methodological core and the finding that matters, not an abstract paraphrase. Records
marked \textsc{metadata only} were surfaced without an abstract in Crossref, Europe PMC or
Semantic Scholar and are listed, not summarised.\par}

\band{Sensing, forecasting \& instrumentation\hfill 3 records}{Amber}

\paperentry{Li et al.\ 2026}
{Spatiotemporally synchronised dual-wavelength photoacoustic--scattering measurement of aerosol optical properties}
{Photoacoustics}
{A sphere--tube coupled resonator that removes the transport, response-time and
microenvironment mismatch inherent to split absorption/scattering instruments. Two
integrating spheres act simultaneously as scattered-light collection cavities and as
acoustic buffers forming a low-frequency \textbf{differential Helmholtz} resonator, while
the central connecting tube supports a high-frequency \textbf{second-order longitudinal}
mode; 450 and 532\,nm lasers are modulated at \textbf{2572} and \textbf{169\,Hz} and
demodulated by an FPGA digital lock-in, so both wavelengths and both quantities are
recovered from one sampling path. Absorption channels calibrated against NO$_2$,
scattering against N$_2$/CO$_2$ Rayleigh scattering; at 1\,s integration and 3$\sigma$
the 532\,nm limits of detection are \textbf{1.2\,Mm$^{-1}$} (absorption) and
\textbf{1.32\,Mm$^{-1}$} (scattering). The validation is the weak part: water mist,
controlled combustion and a field check establish channel isolation and feasibility but
not accuracy against a reference absorption method, and no Allan-variance drift figure is
given, so the 1\,s LOD does not tell you the useful averaging time. For distributed
sensing the transferable idea is architectural rather than economic --- this is a
laboratory instrument, but co-locating absorption and scattering in a shared volume is
what would let a node report single-scattering albedo instead of a mass proxy.}
{10.1016/j.pacs.2026.100864}

\paperentry{Laffolley et al.\ 2026}
{A 3D-printed two-stage virtual impactor ($\mu$SPLIT) for radioactive aerosol size classification and direct analysis}
{ACS Omega}
{Stereolithography-printed virtual impactor separating aerosol into \textbf{$>$10,
1--10 and $<$1\,$\mu$m} aerodynamic classes onto integrated filters for direct
post-collection $\alpha$ counting, designed for Fukushima Daiichi decommissioning where
the instrument must be cheap enough to incinerate rather than decontaminate. Flow path
optimised by CFD with Lagrangian particle tracking; predicted cutoffs \textbf{9.0} and
\textbf{1.3\,$\mu$m}. The paper is unusually candid about its own failure modes:
simulations identify small-particle contamination in the second-stage minor flow and
large-particle contamination in the major flow, and dimensional metrology on the printed
parts found channel contraction and corner rounding --- i.e.\ the as-built geometry is
not the as-designed geometry, which for an impactor moves the cutoff directly.
Experimental validation is a proof of concept only: incense smoke, and Rn-progeny-bearing
NaCl at 1.70\,$\mu$m GMAD putting peak $\alpha$ activity on the middle filter. No
standardised-aerosol calibration and no quantitative wall-loss budget, both of which the
authors name as required next steps. Relevant well beyond radiological work --- additive
manufacture of size-selective inlets is the cheapest available route to giving a
low-cost node a real aerodynamic size cut instead of an optical inference.}
{10.1021/acsomega.6c04521}

\paperentry{Shirkhani et al.\ 2026}
{Design and development of a high-throughput multipattern aerosol deposition device (MARTHA) for respiratory toxicity and health-risk assessment}
{J Aerosol Sci}
{\textsc{metadata only} --- no abstract in Crossref, Europe PMC or Semantic Scholar at the
time of this issue (S2 returns 404 on the DOI). Listed rather than summarised. It is
flagged as the day's highest-priority re-check for a specific reason: the title states an
aerosol-delivery instrument built for respiratory toxicity assessment, which is the exact
instrument\,$\times$\,health-endpoint coupling that the 2--8 August weekly identified as
structurally absent across 146 records. Nothing beyond the title and the author list
(Shirkhani, Salana, Amanatidis, Washeem, Eiguren-Fernandez, Spielman, Nannu Shankar,
Lewis) supports that reading yet.}
{10.1016/j.jaerosci.2026.106882}

\band{Exposure assessment \& modelling\hfill 3 records}{Clay}

\paperentry{Adeniji et al.\ 2026}
{Long-term air quality and meteorological coupling dynamics in Abeokuta, Nigeria}
{PLOS Glob Public Health}
{Twelve years of hourly PM$_{2.5}$, PM$_{10}$, NO$_2$, SO$_2$, CO and O$_3$ (2012--2023)
against six meteorological variables at a single West African site. Particulate matter
dominates and \textbf{consistently exceeds WHO guideline values}, peaking in the
December--March Harmattan; regression against temperature gives \textbf{R$^2$
0.48--0.52}, with negative associations for rainfall and humidity consistent with thermal
stagnation and wet scavenging. The aggregate AQI stays in the ``moderate'' band (39--51)
while the health-risk-adjusted AQI exceeds \textbf{90} seasonally --- the paper's real
contribution, since it demonstrates that an index which averages across pollutants
conceals a PM-driven risk that a health-weighted index exposes. Two load-bearing
limitations: the data provenance for a 12-year hourly record in Abeokuta is not
established in the abstract, and reanalysis or satellite-derived series are common
substitutes at such sites; and single-site, single-city inference cannot separate the
meteorological coupling from local source seasonality (biomass burning peaks in the same
months as the Harmattan). Sub-Saharan Africa remains the most under-monitored region in
this watch and this is the first record from it in six issues.}
{10.1371/journal.pgph.0006683}

\paperentry{Ray et al.\ 2026}
{Evaluating CCN activity through size-resolved hygroscopicity: synergistic observations and simulations in a tropical montane region}
{Atmos Environ}
{\textsc{metadata only} --- no abstract in Crossref, Europe PMC or Semantic Scholar.
Listed, not summarised. Size-resolved hygroscopicity ($\kappa$) is the physical quantity
underneath every relative-humidity correction applied to an optical PM sensor, and
tropical montane conditions are the high-RH regime where those corrections fail worst, so
this is a priority re-check even though the paper's own framing is cloud-condensation
nuclei rather than air quality.}
{10.1016/j.atmosenv.2026.122284}

\paperentry{de la Torre et al.\ 2026}
{A 15-year atmospheric record of intentionally produced persistent organic pollutants in Spain}
{J Hazard Mater}
{Polyurethane-foam passive samplers at \textbf{23 sites, 2009--2023}, evaluating
Stockholm Convention effectiveness. Hexachlorobenzene dominates at a median
\textbf{41\,\%} of $\sum$POPs with a remote-enriched spatial pattern; legacy pesticides
and indicator PCBs show temperature-dependent secondary volatilisation from contaminated
soils by Clausius--Clapeyron analysis, and an \emph{urban} enrichment plausibly driven by
heat-island effects; PBDEs show a \textbf{35-fold} urban-to-remote gradient from
late-regulated BDE-209. Early-banned compounds decline up to \textbf{$-$33\,\%/yr} while
PeCB, iPCBs and BDE-209 are \textbf{stalled or rising to $+$20\,\%}. This is semi-volatile
organic chemistry, not particulate mass, and enters on method rather than substance: a
passive sampler is the cheapest long-record instrument in atmospheric science and this is
a clean demonstration of what a 15-year, 23-site, zero-power network buys --- an
altitudinal profile on Tenerife separating local secondary emission from free-troposphere
clearance. The tradeoff is equally clear: no time resolution, no size information, and
partitioning between gas and particle phases is unresolved throughout.}
{10.1016/j.jhazmat.2026.143155}

\band{Occupational \& indoor exposure\hfill 2 records}{Clay}

\paperentry{Salcines et al.\ 2026}
{Monitoring and risk assessment of airborne graphene oxide released during 3D printing with polylactic acid filaments}
{J Environ Health Sci Eng}
{Tiered exposure assessment to EN 17058:2018 and OECD guidance around fused-deposition
printing of graphene-oxide-loaded PLA, combining \textbf{near-field and far-field particle
counters} with SEM and Raman speciation plus VOC measurement against Spanish and Finnish
2022 occupational limits, with ARIMA used to test whether the observed excursions are
distinguishable from background drift. Nanoparticles and agglomerates were present at low
average number concentration and \textbf{GO was detected in the respirable fraction},
below proposed limits but demonstrably released --- the finding is qualitative
(release occurred) rather than quantitative (how much). Weaknesses are structural: a
pilot with no stated replicate count, number-concentration monitors that cannot speciate,
so the GO attribution rests entirely on offline SEM/Raman of collected material, and
occupational limits for GO that do not formally exist. The methodological point
transfers directly --- paired near-field/far-field counters with an offline speciation
leg is the minimum viable design for attributing a number-concentration excursion to a
source, and it is the design a distributed indoor network would need.}
{10.1007/s40201-026-00998-w}

\paperentry{Shadmani et al.\ 2026}
{A comprehensive health profile and exposure-risk analysis in an iron ore mining workforce}
{J Environ Health Sci Eng}
{Cross-sectional study of \textbf{1,193} Iranian iron-ore mine employees, stratified into
high, moderate and low exposure \emph{by job role}, with physical examination, clinical
chemistry and haematology compared by ANOVA and $\chi^2$ and modelled by linear and
logistic regression. The workforce is young (mean \textbf{34.9\,$\pm$\,7.6\,y}) with low
chronic-disease prevalence; the burden is localised --- hearing \textbf{23.8\,\%}, vision
\textbf{17.8\,\%}, dermatological \textbf{31.3\,\%} --- while most haematological and
metabolic markers did not differ by exposure group. Platelet count, LDL and blood
pressure were \emph{lower} in the high-exposure group, which the authors do not resolve
and which has an obvious competing explanation the paper does not exclude: the
healthy-worker survivor effect plus role-based confounding by age and seniority. The
central methodological weakness for this watch is that \textbf{exposure is a job-role
proxy with no measured dust concentration at all} --- no gravimetric respirable sampling,
no personal monitoring, no size fractionation. A three-level ordinal surrogate cannot
support an exposure--response claim, and the paper's own conclusion (localised sensory and
dermatological effects, not systemic metabolic ones) is a statement about a young cohort
with short cumulative exposure rather than about the hazard.}
{10.1007/s40201-026-01000-3}

\band{Mechanistic toxicology\hfill 1 record}{Deep}

\paperentry{Wei et al.\ 2026}
{Seasonal source shifts drive oxidative potential of size-resolved particulate matter in coastal southeastern China}
{J Hazard Mater}
{\textsc{metadata only} --- no abstract in Crossref, Europe PMC or Semantic Scholar.
Listed, not summarised. Noted because size-resolved oxidative potential is one of the few
routinely measured quantities that is both a particle property and a plausible biological
effect metric, and because a source-driven \emph{seasonal} shift in OP at fixed mass
concentration would be direct evidence against mass as a sufficient exposure metric ---
which is the substantive question behind this watch's instrumentation axis. Whether the
paper supports that reading is unknown until an abstract appears.}
{10.1016/j.jhazmat.2026.143237}

\band{Burden, policy \& mitigation\hfill 2 records}{Coral}

\paperentry{Urso et al.\ 2026}
{Burden of dementia in Europe, 1990--2023: a systematic analysis from the Global Burden of Disease Study 2023}
{Lancet Reg Health Eur}
{GBD 2023 secondary analysis for the EU-27 and the WHO European Region, with non-fatal
outcomes from DisMod-MR 2.1 and dementia-attributable mortality from an excess-mortality
framework. \textbf{7.76\,M (95\,\% UI 6.69--8.72)} prevalent cases in the EU-27 and
\textbf{12.28\,M (10.44--13.90)} across the WHO European Region, both roughly
\textbf{+90\,\% since 1990} despite modest declines in age-standardised prevalence;
dementia moved from the eighth to the \textbf{third} leading cause of death in the EU-27.
\textbf{41\,\% (24.4--57.4)} of EU-27 dementia DALYs are attributed to six modifiable risk
factors, with \textbf{ambient particulate matter pollution named first}, ahead of high
fasting plasma glucose and high BMI, and concentrated in socioeconomically disadvantaged
populations. The limitation is the one every GBD risk-factor attribution carries and this
paper does not resolve: the PM$_{2.5}$--dementia relative risk is imported from a
comparative-risk-assessment function fitted to a thin and heterogeneous cohort literature,
and the \textbf{24.4--57.4} interval is a modelling uncertainty, not a sampling one --- it
does not propagate uncertainty in the exposure surface, which for a 1990--2023 European
back-cast is substantial. Treated here as a burden estimate, not as neurological evidence.
The interval is deliberately absent from the forest plot: an attributable fraction has no
position on a ratio axis.}
{10.1016/j.lanepe.2026.101796}

\paperentry{Sathasivam \& Veluchamy 2026}
{Visual characteristics of PM from \textit{Halymenia dilatata} microalgae biodiesel compared to low-sulfur diesel}
{Environ Sci Pollut Res}
{Single-cylinder Kirloskar CI engine, four microalgae biodiesel blends (B10--B25) against
S10 low-sulfur diesel, PM collected on a textile element in a deceleration chamber,
quantified gravimetrically and characterised by SEM for morphology, deposition and
agglomeration. \textbf{PM mass reduction 10.44--69.79\,\%} across blends. The result does
not vary monotonically with blend ratio --- SEM showed heavier deposition for S10, B15 and
B25 but lighter for B10 and B20 --- which the paper reports but does not explain, and a
non-monotonic dose--response across four points with no stated replication is more likely
measurement noise than combustion chemistry. Three further limits matter: a textile
collection element in a deceleration chamber is not a dilution tunnel and has no
characterised collection efficiency as a function of size, so the gravimetric numbers are
not comparable with any regulated emissions measurement; no size distribution or number
concentration was measured, so a mass reduction is fully compatible with an increase in
ultrafine number; and no engine load or speed conditions are given in the abstract.
Included as a tier-C source-characterisation record.}
{10.1007/s11356-026-38134-8}

\band{Respiratory \& allergic\hfill 1 record}{Sage}

\paperentry{Ni et al.\ 2026}
{Socioeconomic and regional influences on eosinophilic inflammation in chronic rhinosinusitis with nasal polyps}
{World Allergy Organ J}
{Multicentre retrospective series, \textbf{1,047} CRSwNP patients undergoing endoscopic
sinus surgery at three tertiary hospitals in Jiangsu, with paired blood and nasal-mucosa
eosinophil counts at 2008, 2012 and 2016, analysed by correlation, ROC, linear regression
and linear mixed-effects models. Blood eosinophils are a poor surrogate for tissue
eosinophilia (\textbf{$\rho=0.28$}, \textbf{AUC 0.607}, adjusted \textbf{R$^2$ 0.053}) and
vary heterogeneously by hospital and year. City-level GDP correlated with mean blood
($\rho=0.80$, $p=0.014$) and tissue ($\rho=0.72$, $p=0.037$) eosinophil counts, whereas
\textbf{annual mean PM$_{2.5}$ was associated with neither}. The null is the reason this
record is carried, and it should be read as a measurement failure rather than as evidence
of no effect: the exposure is a \emph{city-level annual mean} assigned to surgical
patients at three time points, which is an ecological contrast with essentially no
within-city variance, over an interval (2008--2016) in which Chinese PM$_{2.5}$ fell
steeply and uniformly --- so exposure and calendar time are near-collinear by
construction, and any true association is absorbed by the year term. With eight city-level
observations, the GDP correlations are also fragile. Tier C.}
{10.1016/j.waojou.2026.101440}

\clearpage

% ==================================================================== TABLE
\section*{\sffamily\bfseries\color{Deep}Corpus register}
\vspace{-1.5mm}{\color{Amber}\rule{\linewidth}{1.1pt}}\vspace{2.5mm}

{\fontsize{6.7}{8.4}\selectfont
\setlength{\tabcolsep}{2.4pt}
\begin{longtable}{@{}>{\RaggedRight}p{25mm}>{\RaggedRight}p{41mm}>{\RaggedRight}p{27mm}%
>{\RaggedRight}p{23mm}>{\RaggedRight}p{18mm}>{\RaggedRight}p{15mm}@{}}
\toprule
\textsf{\textbf{First author}} & \textsf{\textbf{Focus}} & \textsf{\textbf{Journal}} &
\textsf{\textbf{Design}} & \textsf{\textbf{Metric}} & \textsf{\textbf{Region}}\\
\midrule
\endfirsthead
\toprule
\textsf{\textbf{First author}} & \textsf{\textbf{Focus}} & \textsf{\textbf{Journal}} &
\textsf{\textbf{Design}} & \textsf{\textbf{Metric}} & \textsf{\textbf{Region}}\\
\midrule
\endhead
\multicolumn{6}{@{}l@{}}{\cellcolor{Mist}\textsf{\textbf{\textcolor{Deep}{Neuro / mental health}}}}\\[0.6mm]
\href{https://doi.org/10.1001/jamanetworkopen.2026.25015}{Dales} & Wildfire-sourced PM$_{2.5}$ and ED visits for migraine and other primary headache & \textit{JAMA Netw Open} & Case-crossover & PM$_{2.5}$ (wildfire) & Canada\\
\href{https://doi.org/10.1093/ageing/afag215}{Glover} & Umbrella meta-review of urban environmental exposures and cognitive health & \textit{Age Ageing} & Umbrella review / meta-analysis & PM$_{2.5}$ & Global\\
\href{https://doi.org/10.1016/j.nsa.2026.107019}{Goossens} & In utero air pollution exposure and pre-/postnatal brain development & \textit{Neurosci Appl} & Narrative review & PM$_{2.5}$ / PM$_{10}$ / PM$_{\mathrm{coarse}}$ & Global\\
\href{https://doi.org/10.1016/j.isci.2026.116617}{Ma} & PM$_{2.5}$-induced neurotoxicity and depression: exposure risks and mechanisms & \textit{iScience} & Narrative review & PM$_{2.5}$ + components & Global\\
\href{https://doi.org/10.3760/cma.j.cn112150-20250804-00752}{Pan} & Maternal PM exposure and offspring neurodevelopment: effects and mechanisms & \textit{Zhonghua Yu Fang Yi Xue Za Zhi} & Narrative review & PM (bulk) & China\\
\addlinespace[0.9mm]
\multicolumn{6}{@{}l@{}}{\cellcolor{Mist}\textsf{\textbf{\textcolor{Teal}{Cardiovascular \& metabolic}}}}\\[0.6mm]
\href{https://doi.org/10.3760/cma.j.cn112338-20251120-00843}{Huang} & PM$_{2.5}$ and dyslipidaemia in rural Chinese children; dietary effect modification & \textit{Zhonghua Liu Xing Bing Xue Za Zhi} & Cross-sectional (multistage) & PM$_{2.5}$ & China\\
\href{https://doi.org/10.32394/pe/218966}{Kazmierska} & Air pollution and cardiovascular disease: review of current data & \textit{Przegl Epidemiol} & Narrative review & PM$_{2.5}$ / PM$_{10}$ & Global\\
\href{https://doi.org/10.1007/s00125-026-06814-2}{Xi} & Air pollution, built environment and temperature vs. prediabetes phenotypes & \textit{Diabetologia} & Prospective cohort & PM$_{2.5}$ & Germany\\
\href{https://doi.org/10.1161/JAHA.125.046917}{Yang} & Short-term PM$_{2.5}$ and hospitalisation/ED use in atrial fibrillation & \textit{J Am Heart Assoc} & Case-crossover & PM$_{2.5}$ (wildfire-influenced) & USA\\
\addlinespace[0.9mm]
\multicolumn{6}{@{}l@{}}{\cellcolor{Mist}\textsf{\textbf{\textcolor{Sky}{Reproductive \& developmental}}}}\\[0.6mm]
\href{https://doi.org/10.1016/j.jes.2025.10.043}{Chen} & Prenatal air pollution, fetal growth trajectory phenotypes and child neurodevelopment & \textit{J Environ Sci (China)} & Prospective cohort & PM$_{2.5}$ / PM$_{10}$ & China\\
\href{https://doi.org/10.1016/j.ecoenv.2026.120566}{Huang} & Joint prenatal bisphenol + air-pollutant mixture exposure and preschool IQ & \textit{Ecotoxicol Environ Saf} & Prospective cohort (mixtures) & PM$_{2.5}$ / PM$_{10}$ & China\\
\href{https://doi.org/10.1007/s40572-026-00552-8}{Peterson} & Ambient air pollution, PAHs and pubertal development: critical review & \textit{Curr Environ Health Rep} & Critical review & PM$_{2.5}$ components & Global\\
\href{https://doi.org/10.1016/j.envres.2026.125313}{Xie} & Long-term PM$_{2.5}$ and timing of pubertal development in boys and girls & \textit{Environ Res} & Prospective cohort & PM$_{2.5}$ & China\\
\addlinespace[0.9mm]
\multicolumn{6}{@{}l@{}}{\cellcolor{Mist}\textsf{\textbf{\textcolor{Sage}{Respiratory \& allergic}}}}\\[0.6mm]
\href{https://doi.org/10.1007/s00484-026-03289-z}{Bayar Muluk} & Climate change and otolaryngology: airway disease impacts and emission reduction & \textit{Int J Biometeorol} & Narrative review & PM$_{2.5}$ & Global\\
\href{https://doi.org/10.4103/mgr.MEDGASRES-D-25-00164}{Liao} & Short-term air pollution, respiratory symptoms and lung function in asthma & \textit{Med Gas Res} & Panel study & PM$_{2.5}$ / PM$_{10}$ & China\\
\href{https://doi.org/10.3390/epidemiologia7040097}{Sanchez de Prada} & Ambient PM and asthma-coded admissions in children 0-3 y, Spain & \textit{Epidemiologia} & Ecological panel & PM$_{2.5}$ / PM$_{10}$ & Spain\\
\href{https://doi.org/10.1016/j.jhazmat.2026.143076}{Yang} & Allergenic vs. toxic risk of airborne seed hairs compared with ambient PM$_{2.5}$ & \textit{J Hazard Mater} & Field measurement & PM$_{2.5}$ + bioaerosol & China\\
\addlinespace[0.9mm]
\multicolumn{6}{@{}l@{}}{\cellcolor{Mist}\textsf{\textbf{\textcolor{Violet}{Mechanistic toxicology}}}}\\[0.6mm]
\href{https://doi.org/10.1016/j.jes.2025.12.074}{Ding} & PM$_{2.5}$ promotes vascular calcification via miR-27b-5p/HES1 signalling & \textit{J Environ Sci (China)} & In vitro + organoid + murine & PM$_{2.5}$ & China\\
\href{https://doi.org/10.1007/s00360-026-01698-5}{Maraschi} & Osmoregulatory disruption and metal bioaccumulation from settleable atmospheric PM & \textit{J Comp Physiol B} & Controlled exposure (ecotox) & Settleable APM & Brazil\\
\href{https://doi.org/10.1016/j.jes.2025.11.018}{Wu} & Macrophage NLRP3 inflammasome in real-ambient PM-induced hepatic glucose dysregulation & \textit{J Environ Sci (China)} & Real-ambient murine + co-culture & Real-ambient PM & China\\
\addlinespace[0.9mm]
\multicolumn{6}{@{}l@{}}{\cellcolor{Mist}\textsf{\textbf{\textcolor{Clay}{Occupational \& indoor}}}}\\[0.6mm]
\href{https://doi.org/10.3389/fmedt.2026.1796828}{Chen} & Nested suction instrument to cut intraoperative surgical-smoke PM$_{2.5}$ & \textit{Front Med Technol} & Bench / device experiment & PM$_{2.5}$ (surgical smoke) & China\\
\href{https://doi.org/10.1093/annweh/wxag055}{Cole} & Determinants of respirable crystalline silica exposure in Australian tunnelling & \textit{Ann Work Expo Health} & Observational exposure study & Respirable dust / RCS & Australia\\
\href{https://doi.org/10.1016/j.jhazmat.2026.143081}{Ding} & Stage-resolved hair Zn-isotope dynamics as a firefighting combustion-exposure tracer & \textit{J Hazard Mater} & Repeated-measures biomarker & Combustion-derived PM & China\\
\href{https://doi.org/10.3390/jox16040119}{Ramsden} & Xenobiotic hazards in aircraft cabin air & \textit{J Xenobiot} & Review + risk economics & Ultrafine particles & Global\\
\addlinespace[0.9mm]
\multicolumn{6}{@{}l@{}}{\cellcolor{Mist}\textsf{\textbf{\textcolor{Amber}{Sensing, forecasting \& instrumentation}}}}\\[0.6mm]
\href{https://doi.org/10.1016/j.jenvman.2026.130515}{Shi} & Adaptive spatiotemporal graph transformer for multi-site multi-step PM$_{2.5}$ forecasting & \textit{J Environ Manage} & Deep-learning model & PM$_{2.5}$ & China / India\\
\href{https://doi.org/10.3390/bios16070380}{Tsuchiya} & Seasonal and indoor-outdoor characteristics of fluorescent aerosol particles in offices & \textit{Biosensors} & Real-time sensor campaign & Bioaerosol / coarse & Japan\\
\addlinespace[0.9mm]
\multicolumn{6}{@{}l@{}}{\cellcolor{Mist}\textsf{\textbf{\textcolor{Amber}{Exposure assessment \& modelling}}}}\\[0.6mm]
\href{https://doi.org/10.1029/2025GH001621}{Solorzano-Araque} & Multiscale exposure modelling: LOTOS-EUROS + mobility surveys + particle toxicity & \textit{GeoHealth} & CTM + mobility integration & PM$_{2.5}$ / PM$_{10}$ & Colombia\\
\addlinespace[0.9mm]
\multicolumn{6}{@{}l@{}}{\cellcolor{Mist}\textsf{\textbf{\textcolor{Coral}{Burden, policy \& mitigation}}}}\\[0.6mm]
\href{https://doi.org/10.1002/hsr2.72886}{Karamzad} & Ischaemic heart disease burden in MENA 1990-2021; PM attributable fraction 28.2\% & \textit{Health Sci Rep} & GBD secondary analysis & Ambient PM & MENA\\
\href{https://doi.org/10.1371/journal.pone.0347989}{Rees} & HEPA filtration to reduce respiratory infection in care homes & \textit{PLoS One} & RCT process evaluation & Filtered PM & UK\\
\href{https://doi.org/10.1016/j.jes.2026.01.009}{Shan} & Policy-driven emission reductions and post-pandemic rebound, high-resolution inventories & \textit{J Environ Sci (China)} & Emission inventory & PM emissions & China\\
\addlinespace[0.9mm]
\multicolumn{6}{@{}l@{}}{\cellcolor{Mist}\textsf{\textbf{\textcolor{Slate}{Other clinical endpoints}}}}\\[0.6mm]
\href{https://doi.org/10.3390/dj14070440}{Lupita} & PM$_{2.5}$, preventive-care capacity and caries/edentulism burden across the EU & \textit{Dent J} & Ecological panel & PM$_{2.5}$ & EU-27\\
\href{https://doi.org/10.1016/j.jdin.2026.06.008}{Park} & Long-term PM exposure and risk of androgenetic alopecia & \textit{JAAD Int} & Nationwide cohort & PM$_{2.5}$ / PM$_{10}$ & South Korea\\
\href{https://doi.org/10.3171/2026.3.PEDS2615}{Sun} & Environmental exposure and area deprivation vs. paediatric hydrocephalus outcomes & \textit{J Neurosurg Pediatr} & Registry cohort + geospatial & PM$_{2.5}$ (satellite-derived) & USA\\
\addlinespace[0.9mm]
\bottomrule
\end{longtable}}


\vspace{4mm}
\begin{keybox}
{\sffamily\bfseries\footnotesize\color{Deep}Provenance and method}\\[1.2mm]
{\fontsize{7.2}{8.8}\selectfont
Entry window \textbf{2026-08-09 to 2026-08-09}, contiguous with the 8 August issue; no
gap. Legs run: \textbf{PubMed} connector, two independent queries against
\texttt{[Date - Entry]} --- health axis 6 hits, instrumentation axis 4 (overlapping),
plus a third broad sweep returning 15 for the day. \textbf{Four of those 15 were already
shipped in the 4 August issue} and re-surfaced on re-indexing; they were caught by
\texttt{seen.json} on both PMID and DOI. The local harvester returned 5 health and
1 sensing. \textbf{Europe PMC} on \texttt{CREATION\_DATE} returned 6, of which 3 were
already shipped and 1 was trade news; \textbf{2 carried}, including the day's second
instrument paper. \textbf{Crossref by ISSN} across 17 journals returned 9, \textbf{3
carried}. \textbf{Consensus} $\times$3 sweeps returned 20 each, capped at 3 per sweep by
the free tier; all 9 returned DOIs resolved to Crossref creation dates between 12 January
and 15 July 2026 and were rejected as outside the window, 0 carried. \textbf{arXiv}
50 records, none published on or after 1 August. \textbf{OpenAlex} HTTP 429, seventh
consecutive run. \textbf{ClinicalTrials.gov} 0 registrations or updates on both the
condition and the intervention query.
\textbf{16 records rejected}, each with a reason in \texttt{state/rejected.jsonl}.
Three carried records have no abstract in any of the three checked sources and are
listed without a summary; they are marked \textsc{metadata only} in the digest and in
the architecture donut.
\textbf{No forest plot is generated for this issue}: the three interval estimates in the
corpus are a population-attributable fraction and two prevalence counts, none of which is
ratio-scale, and \texttt{f5\_forest.png} is therefore intentionally absent rather than
missing.
Subtopic, design, particle-metric and geography labels are assigned from title, abstract,
keywords and MeSH; assignment is single-label by dominant contribution and is a
judgement, not a controlled vocabulary.
All DOI links resolve to the publisher of record. Content derived from PubMed and
Crossref metadata; abstracts remain the copyright of their respective publishers.\par}
\end{keybox}

\end{document}
